%% sections/03-methodology.tex

\section{Methodology}
\label{sec:methodology}

\subsection{The Promotion Pipeline}

A player style is a named behavioral design pattern: a recurring class of PC behavior
that is (a) enabled by specific character-sheet features, (b) identifiable from
session-log evidence, (c) replicable across campaigns with different party archetypes,
and (d) actionable as a design prescription.

The pipeline that promotes innovation instances to ratified player styles operates as
follows.

\paragraph{Stage 1: Innovation logging.}
A session-gate reviewer reads each completed session log and flags \textit{innovation
instances}: moments when a PC's action or speech exceeds what the current playtest
rubric anticipates. An instance is logged with (a) the session and scene reference,
(b) a description of the behavior, (c) the character-sheet feature that appears to
have enabled it, and (d) a proposed thematic cluster label (e.g., ``sheet-deep-reader
candidate'', ``craft-reception candidate'').

\paragraph{Stage 2: Threshold accumulation.}
A cluster crosses the promotion threshold when it contains 3 or more instances spanning
2 or more distinct sessions. The threshold is intentionally more demanding than the
rubric amendment threshold (which requires only 2 instances in a dimension), because
player styles are claimed to be design-enabled rather than session-specific. The
instances must also be \textit{substantively distinct}: three instances of the same
mechanic in different scenes do not qualify; three instances showing the pattern from
different contexts do.

\paragraph{Stage 3: Player-style proposal.}
When the threshold is crossed, a player-style entry is drafted. The entry format is:
\begin{description}
  \item[Definition.] What the behavior is, in PC-design terms (not player-personality
    terms). The definition must be specific enough to generate design prescriptions.
  \item[Triggers/Signals.] Observable session-log indicators that the style is active:
    specific types of PC action, speech patterns, or narrative moves.
  \item[Design triggers.] Character-sheet features that enable the style; features
    that undermine it.
  \item[Cross-applicability.] Design features in adventures or campaign spines that
    amplify instances.
  \item[Confirmed sub-types.] Named variants that have accumulated 3 or more cross-campaign
    instances of their own.
  \item[Cross-campaign validation status.] Which campaigns have produced instances;
    which campaign archetypes the style has been confirmed in.
\end{description}

\paragraph{Stage 4: Ratification.}
A second review confirms that the promotion-threshold instances are substantively
distinct and that the definition correctly describes the observed pattern. Ratification
does not require additional instances; it requires that the existing instances are
correctly categorized.

\paragraph{Stage 5: Cross-campaign monitoring.}
After ratification, the style is tracked across subsequent campaigns. The relevant
questions are: (a) Does the pattern recur? (b) Does it fire in campaigns with different
motivational structures (grief-themed, duty-based, professional, relational)? (c) Do
consistent sub-variations accumulate to warrant sub-type naming?

A style that has been ratified but has produced no cross-campaign instances is noted
as archetype-specific until evidence accumulates. No style in the corpus reached the
end of the 7-campaign observation window in this state.

\subsection{Threshold Calibration}

The promotion threshold (3 instances across $\geq$2 sessions) is a design choice
that requires empirical justification. We calibrated it retrospectively against the
full 7-campaign corpus by simulating alternative thresholds.

\paragraph{Threshold 2/2 (2 instances, $\geq$2 sessions).}
Under a 2/2 threshold, all 7 promoted styles would have been promoted, but promoted
earlier: mean promotion point would have advanced by 1.8 sessions. However, 3
candidate patterns that did not generalize would also have been promoted under this
threshold. Each of these patterns appeared in 2 sessions but did not recur in
subsequent campaigns, and none produced enough cross-campaign instances to satisfy
the retrospective validity criterion used for the 7 ratified styles. These 3 patterns
represent false positives that the 3/2 threshold correctly excluded.

\paragraph{Threshold 4/2 (4 instances, $\geq$2 sessions).}
Under a 4/2 threshold, the two highest-instance styles --- sheet-deep-reader
($\sim$50 instances) and craft-witness ($\sim$45 instances) --- would still have
been promoted. However, documentary-witness (7 instances) and act-without-announcement
intra-session (8 instances) would not yet have been promoted at the point in the
corpus where the 7-campaign observation window closes. Both styles proved
session-invariant in the campaigns where they appeared; a 4/2 threshold would have
delayed their ratification without producing more reliable promotions.

\paragraph{Conclusion.}
The 3/2 threshold represents a conservative empirical choice: it produced 0 false
promotions (no promoted style was later demoted) and promoted all patterns that proved
session-invariant. We recommend treating it as a starting point for new campaigns,
with recalibration if false promotion rates exceed 10\%. The threshold is not
theoretically derived; it is a pragmatic calibration point supported by retrospective
analysis of this corpus. Campaigns with substantially different session cadences,
party archetypes, or session lengths may warrant lower thresholds (if session count
is compressed) or higher thresholds (if session heterogeneity is high and false
promotions are costly).

\subsection{The Seven Promoted Styles}

Table~\ref{tab:styles-overview} provides a summary of all 7 styles.

\begin{table*}[t]
\centering
\caption{The 7 Promoted Player Styles: Definition, Instance Count, Promotion Point,
         and Confirmed Campaign Range}
\label{tab:styles-overview}
\begin{tabular}{@{}lp{5cm}rp{2.5cm}@{}}
\toprule
Style & Core definition & Instances & Confirmed in \\
\midrule
sheet-deep-reader
  & PC reads sheet as live agency source; surfaces backstory without DM prompting
  & $\sim$50 & C1--C7 (all) \\
craft-witness
  & Named daily practice is an atmospheric antenna; the ritual is the reception
  & $\sim$45 & C1--C5 \\
npc-arc-completion trigger
  & Party behavior meets NPC firing condition; arc fires by character logic, not
    DM direction
  & $\sim$40 & C1--C7 (all) \\
act-without-announcement (cross-session)
  & Post-naming behavioral change without declaration; the act is the completion
  & $\sim$25 & C1--C4 \\
verbal-documentation-arc
  & Oral-practice PC meets vocabulary limit; 3-stage arc to naming-the-thing
  & $\sim$15 & C2--C3 \\
documentary-witness
  & Documentation-as-reception; limit-entry is arc completion
  & $\sim$12 & C4, C7 \\
act-without-announcement (intra-session, v1.12)
  & Same-session naming event and behavioral consequence both present
  & $\sim$8 & C2--C4 \\
\bottomrule
\end{tabular}
\end{table*}

\subsubsection{Sheet-Deep Reader ($\sim$50 instances)}

\textit{Definition.} A PC reads their own character sheet as a live source of agency
throughout play, surfacing grief-paragraphs, prior-adventure hooks, named relationships,
and institutional training as current-scene decision inputs without being prompted by
the DM.

The style has three confirmed sub-types:

\textit{Grief-paragraph-dialogue} ($\sim$18 instances): the specific sensory memory or
named loss from the PC's grief paragraph surfaces as spoken dialogue at a moment of
high emotional weight. The character sheet content becomes the character's exact words.
In Campaign 2, Thessaly said: ``I cast a spell at my Test that I was not supposed to
cast. It cost someone something I can't give back.'' This is verbatim grief-paragraph
content, spoken without DM prompting.

\textit{Commitment-accounting} ($\sim$16 instances): a PC tracks promises made across
sessions and names them unprompted before anyone can forget. In Campaign 2 Session 5,
Calder named the witness-list without being asked: Mira, Deva, Vorn, Sevven, Morreth's
chapel. Five names, in order, none of which the DM had referenced.

\textit{Practitioner-by-material-relationship} ($\sim$16 instances): a PC identifies
other practitioners by how they handle their materials rather than by class or title.
In Campaign 2, Lenne identified three caster-school affiliations from posture,
text-handling, and annotation habits across three consecutive sessions.

\textit{Design trigger.} PC sheets with detailed grief-paragraphs, named losses,
local-origin notes, institutional training, and prior-adventure hooks. Generic
backstory fields (``bond: my village'', ``flaw: stubborn'') undermine the style.

\subsubsection{Craft-Witness ($\sim$45 instances)}

\textit{Definition.} A PC with a named daily craft-ritual or attention-practice
receives atmospheric flags other PCs miss. The craft practice functions as a
reception antenna: not because the PC has a higher Wisdom modifier, but because the
practice gives them a specific, named mode of attention.

This style is \textit{session-invariant} once established for a given PC. Grom Ironhand
(Campaign 1) produced craft-witness instances in every session he appeared in across
all 7 Campaign 1 sessions (16 instances total). Orik Flintback (Campaign 2) produced
craft-witness instances in all 7 Campaign 2 sessions (14 instances total). Neither PC
missed a session.

The style fires from secular professional practice (Orik's hall-guard structural
assessment: ``door is set wrong, frame doesn't carry the load this passage suggests'')
as well as religious practice (Grom's forge-priest hearth-greeting: ``greet hearths:
thou'rt seen''). Practice type is irrelevant; the specificity of the named practice
is what enables the style.

\textit{Design trigger.} PC sheets that specify a named daily ritual with a specific
physical gesture or observation habit. Generic entries (``meditates each morning'') do
not produce instances. Specific entries (``greets every forge and hearth aloud'' or
``assesses structural load-paths when entering any built space'') do.

\subsubsection{NPC-Arc-Completion Trigger ($\sim$40 instances)}

\textit{Definition.} Party behavior meets the firing condition specified in an NPC's
arc-completion design; the NPC fires the arc by character logic rather than DM
direction. This style is unusual in that the enabling design work is on the NPC side,
not the PC side: it requires NPCs designed with explicit arc-completion text and
specific firing conditions.

This is the only style in which the PC's behavioral contribution is diffuse: the PCs
do not exhibit a distinctive act; they provide the conditions in which the NPC's arc
fires. The style was identified when reviewers noticed that NPC arc-completions were
consistently triggered by specific party behaviors (patience, a named act of witness,
a question asked without advocacy) rather than by DM narration.

\textit{Design trigger.} NPCs with written arc-completion text specifying what the NPC
will say or do and under what conditions they will say or do it. Without this design
work on the NPC side, the trigger-side pattern cannot fire.

\subsubsection{Act-Without-Announcement, Cross-Session ($\sim$25 instances)}

\textit{Definition.} A PC acts differently in a later session as a direct result of a
prior naming event, without declaration. The naming event is session N (a grief named,
a longing named, a commitment named); the behavioral change is session N+1 or later.
The behavior is the completion; the PC does not explain the change.

Canonical instance: Thera Voss (Campaign 1) named her longing for Adda in Session 6
(``I counted too''). In Session 7, on the road home from Vaenshold, she cut bread
long-ways --- the first time since Adda died when she was eight. She was aware of
doing it. She did not explain. The act follows the naming; the naming is not
repeated.

Sera Ashfall (Campaign 2) produced 11 consecutive instances across all 7 Campaign 2
sessions, making this the session-invariant instance of the cross-session form: once
a PC's act-without-announcement pattern is established, it fires every session.

\textit{Design trigger.} A prior-session naming event encoded in the PC sheet (either
as grief-paragraph content or as a between-session note); a subsequent scene in which
the named thing is available to act from.

\subsubsection{Verbal Documentation-Arc ($\sim$15 instances)}

\textit{Definition.} A PC whose primary documentation mode is oral encounters content
that exceeds their current vocabulary. The arc runs: Stage 1 (cannot-categorize:
the oral system breaks down); Stage 2 (gap-identified: the PC names what their
framework cannot hold); Stage 3 (names-the-thing: oral transmission completes,
transforming the content in the speaking).

Source: Calla Fosse (Campaign 3). Stage 1: Insight fumble (total 6, DC 12; framework
broken, Calla says the account ``comes out sideways''). Stage 2: Religion 14 (identifies
the correct sequence; the framework is reshaped). Stage 3: at the Keystone, Roleplaying
23 (campaign high). Calla did not recite what she had been holding; she spoke
Velantha's true account into existence. The panel review noted that the account changed
in the speaking.

Documentary sub-type (Lenne, Campaign 2): written variant. Stage 1: documentation
suspended (cannot file the observation). Stage 2: vocabulary found (the accountability
clause draft). Stage 3: Lenne's clause language spoken verbatim at the compact speaking.

\textit{Design trigger.} PC sheet with oral documentation listed as the primary
practice heuristic; a scene designed to exceed the PC's current vocabulary.

\subsubsection{Documentary-Witness ($\sim$12 instances)}

\textit{Definition.} A PC whose primary atmospheric reception mode is external
documentation accumulates a cross-session record. The arc completes when the PC
produces a limit-entry: an observation that cannot be filed into any existing category,
signaling that the documentation itself has become inadequate.

Source: Caelith Vor (Campaign 4). Six sessions of grey-blue map annotations, each
adding a convergence point. Session 7: after a failed Insight roll, Caelith wrote
``unable to categorize'' for the first time. That notation is the arc's completion:
the six-session silent-reception chain was the reception; the limit-entry was its
conclusion.

\textit{Design trigger.} Written documentation practice listed as the primary heuristic;
a cross-session information thread with a limit-entry structurally available.

\subsubsection{Act-Without-Announcement, Intra-Session (v1.12, $\sim$8 instances)}

\textit{Definition.} The naming event and its behavioral consequence both occur in the
same session. The intra-session form was ratified as a confirmed sub-type in rubric
amendment v1.12 after three instances across two campaigns showed the same structure
as the cross-session form but compressed into a single session.

Canonical instance: Sava (Campaign 4) received a token in Scene 2 and placed it
beside her suture kit in Scene 5 without showing the party. The token placement was
not commented on; no character narrated it. It followed the naming-moment (the token
receiving) by three scenes.

\textit{Design trigger.} A naming-moment earlier in the session and a scene later in
the same session where the named thing is available to act from. The PC sheet need not
specify this explicitly; the pattern emerges when the session structure provides both.

\subsection{PC-Authored Finale Deliverables}

Across 4 campaigns, a PC-authored deliverable --- a speech, written document, forgery,
or deliberate act --- was constructed across multiple sessions and deployed at the
campaign finale. Table~\ref{tab:deliverables} summarizes these deliverables.

\begin{table*}[t]
\centering
\caption{PC-Authored Finale Deliverables: Campaign, PC, Type, Build Sessions, Key Sheet Feature}
\label{tab:deliverables}
\begin{tabular}{@{}llllrp{3.5cm}@{}}
\toprule
Campaign & PC & Type & Build sessions & Words & Key sheet feature \\
\midrule
C1 Moon-Silver & Grom Ironhand & Spoken confession & S03--S07 (5 sessions) & $\sim$280
  & Reorx-craft heuristic + Duren grief paragraph \\
C2 Conclave Compact & Thessaly Nightmantle & Compact opening sentence & C2-S01--C2-S07 (7 sessions) & $\sim$45
  & Named Test wound; decision-order heuristic \\
C4 Thorngate Watch & Davan Fostermark & Accounting speech & C4-S04--C4-S07 (4 sessions) & $\sim$480
  & Honor-accounting heuristic; Orel grief paragraph \\
C5 Thieves Guild & Mira Calloway & False dossier (MARTEN-7) & C5-S01--C5-S07 (7 sessions) & $\sim$120
  & Framework-accumulation heuristic; professional code \\
\bottomrule
\end{tabular}
\end{table*}

All four deliverables share a construction pattern: each begins as a structural
invitation in the PC sheet (a grief paragraph, a professional code, a decision-order
heuristic) and accumulates substance across sessions as campaign events provide
material. The deliverable is not scripted by the designer; it is grown by the system
playing the PC in response to campaign events. The ending feels like an inevitable
conclusion because the PC's sheet and the campaign's events accumulated toward it
together across multiple sessions.

Grom's Silver-Ingot Confession named two shames at the forge: Aelwen's (things made
to hold grief; the withheld tool) and his own (``I was slow and I was proud and I
called it patience''). Neither was scripted; both were shaped by session events acting
on the PC sheet's Reorx-forge-craft identity across five sessions.

Thessaly's compact opening sentence (``We speak this not to bind the Orders to each
other but to bind each Order to the people who live beneath the consequences of what
the Orders do'') was written in the reconstruction-document margin during Campaign 2
Session 6 and spoken aloud at the Session 7 finale. It compressed seven sessions of
wound-naming, witness-gathering, and compact-clause recovery into a single sentence.

Mira Calloway's false dossier (MARTEN-7) accumulated across 7 sessions as her
framework-accumulation heuristic gathered six components. At the finale, she deployed
the completed dossier to neutralize the campaign's antagonist through professional
action rather than confrontation: ``Pale described party as low-yield fencing network,
no intelligence activity, RI value: none, recommend inactive.'' This is a Route B
outcome consistent with the sheet's professional-code identity.
